\section{Leanコミュニティとリーマン予想の形式化}
AnthropicのLevent AlpögeがClaudeと共同で作成したとされるリーマン予想のLean形式化について、Lean／Mathlibコミュニティが正しさに強い確信を示したとの報告が議論を集めた。Grokの収集内容によれば、約1700万行の形式化とGitHub上での確認作業が言及されている。関連するコミュニティ声明は11月21日頃の観測窓外資料であり、今回の窓内ではそれを参照する二次的な投稿が中心だった。Claudeの寄与の程度についても投稿上の主張として扱う。今後は数学的解説や独立した検証の進展が焦点となる。
\dailyxnote{観測窓内の公式一次X投稿は未特定。声明・リポジトリは窓外資料として言及。}

\section{ObsidianとGrok Botの公式連携}
Obsidian公式アカウントは、Grok BotがObsidianのvaultへ接続し、ノートの作成、検索、管理を行えるプラグインを発表した。情報整理環境に対話型エージェントを直接接続する取り組みとして関心が続いている。公式X投稿は11月19日頃で観測窓外だが、窓内でも関連する言及が見られた。利用時にはvaultへのアクセス権限や操作の承認、同期範囲の確認が重要となる。
\dailyxsource{Obsidian公式（観測窓外の発表）}{https://x.com/obsdmd/status/2123484964268413341}

\section{Hermes Agentの承認優先型タスク提案}
Hermes Agentで、プロアクティブなタスク提案を人間の承認に従わせる変更が議論されているとGrokが報告した。承認カードを介した同意を必要とする仕組みが言及され、エージェントが自律的に動く範囲と人間の判断をどう分けるかが論点となる。ただし開発者の窓内一次投稿や公式PRは特定できておらず、安定版が0.23.3のままとする記述も未検証である。現時点では実装確定ではなく関連投稿で観測された議論として扱う。

\section{PerplexityにおけるGemini 4 Neon採用情報}
PerplexityのCEOがTeammatesの一部をGemini 4 Neonへ移行したと投稿した、との二次情報が共有された。複数モデルを用途ごとに選ぶサービスでの採用事例として注目されるが、Grokは当該CEO投稿を観測窓内で特定できていない。移行範囲や性能比較、正式な提供条件は確認されておらず、採用が確定した事実としては扱わない。今後の公式説明が待たれる。

\section{Cloudflare Agents SDK 0.34の待機機構}
Cloudflare Agents SDK v0.34.0では、人間の承認や期限を待つ間、Durable Objectを稼働状態に維持し続けないための機能が追加されたとGrokが収集した。承認ゲートやサブエージェントの待機を伴う長時間ワークフローに関係する変更であり、実行資源と待機状態を分離する設計が注目される。Pi 1.6対応も言及された。ただし根拠となる公式changelogは観測窓外の更新であり、窓内では関連する議論が中心である。

\section{Inkling 2 Miniのローカル推論テスト}
Thinking MachinesのInkling 2 Miniについて、Strataエンジンを用いたローカルテストの結果が投稿された。投稿者は生成速度の中央値を53.3トークン毎秒とし、HTTP 503エラーが約4割発生したと報告した。NaNに関する不具合はエンジン側に起因する可能性があるという見方も示された。これらは単一投稿者による未検証の測定・推定であり、モデル固有の性能や信頼性を示す一般的な結果ではない。モデル自体の発表は観測窓外だった。
\dailyxsource{サクトアスラのローカルテスト投稿}{https://x.com/sakusaku_japan/status/2124625047818220018}

\section{OpenAIのPerplexity曲線に関する投稿}
OpenAIのchief scientistによるとされるエッセイ「Perplexity and Intelligence」を巡り、log-perplexityが0.68付近に達すると次の能力上昇が起こるという予測が二次的に紹介された。投稿ではリーマン予想の形式検証を巡る話題とも関連づけられたが、Grokはエッセイの一次資料を確認できていない。数値や予測は投稿者による要約・解釈として扱い、OpenAIの正式見解や実証済みの閾値とは区別する。
\dailyxsource{Muscle AIによる二次投稿}{https://x.com/themuscleai/status/2124625688426856699}

\section{AnthropicのIPO後の市場反応への言及}
AnthropicのIPO後の株価変動や市場の受け止め方を巡る二次投稿が観測された。ただしGrokの収集では具体的な価格や変動率、公式IR情報は確認されていない。生成AI企業への市場評価が引き続き関心を集めるという観測にとどまり、個別の株価推移や上場条件を確定情報として記載しない。公式な開示資料が今後の確認先となる。
